%-------------------------------------
% This template comes from Anish Athalye (Unofficial University of Cambridge Poster Template).
%
% The Poster Template has been modified by Dr. Rahul Raoniar to fulfill
% B.Tech/Master/Ph.D./PostDoc student's poster presentation requirements.
%-------------------------------------

\documentclass[final,20pt]{beamer}

% ====================
% Packages
% ====================

\usepackage[T1]{fontenc}
\usepackage{lmodern}
\usepackage[orientation=portrait,size=a0,scale=1.02]{beamerposter}
\usetheme{gemini}
\usecolortheme{nott}
\usepackage{graphicx}
\usepackage{booktabs}
\usepackage{tikz}
\usepackage{pgfplots}
\pgfplotsset{compat=1.14}
\usepackage{anyfontsize}
\usepackage{xcolor}
\usepackage[skip=2pt,font=large]{subcaption}
\usepackage{adjustbox}
\usepackage{amsmath,amssymb}

% ----------------------------------
% For plotting study methodology
% ----------------------------------

\usetikzlibrary{shapes.geometric, arrows, positioning}

\tikzstyle{startstop} = [rectangle, rounded corners, minimum width=3cm, minimum height=1cm, text centered, text width = 10cm, draw=black, fill=white]
\tikzstyle{process} = [rectangle, minimum width=3cm, minimum height=1cm, text centered, text width = 6cm, draw=black, fill=white, text width = 10cm]
\tikzstyle{arrow} = [ultra thick,->,>=stealth]

% ====================
% Lengths
% ====================

% Two-column A0 portrait layout: 3 separators + 2 columns = paper width.
\newlength{\sepwidth}
\newlength{\colwidth}
\setlength{\sepwidth}{0.025\paperwidth}
\setlength{\colwidth}{0.4625\paperwidth}

\newcommand{\separatorcolumn}{\begin{column}{\sepwidth}\end{column}}

% ====================
% Title
% ====================

\title{Improving Emotional Text-to-Speech in Low-Resource Languages\\via Reinforcement Learning}

\author{Hao-Chun Hsieh \and Pao-Hua Hsu \and Nien-Tse Wu}

\institute[shortinst]{Institute of Artificial Intelligence Innovation, National Yang Ming Chiao Tung University}

% ====================
% Footer (optional)
% ====================

\footercontent{
  \href{mailto:handsome.ii14@nycu.edu.tw}{\textbf{handsome.ii14@nycu.edu.tw}} \hfill
  \textbf{NYCU 535514 Reinforcement Learning} \hfill
  \href{mailto:mattwu.ii14@nycu.edu.tw}{\textbf{mattwu.ii14@nycu.edu.tw}}}

% ====================
% Logo (optional)
% ====================

\logoleft{
\includegraphics[height=7cm]{logos/NYCU_Logo_C.png}}

% ====================
% Body
% ====================

\begin{document}

\begin{frame}[t]
\begin{columns}[t]
\separatorcolumn

\begin{column}{\colwidth}

% ----------------------------------
% Abstract
% ----------------------------------
  \begin{block}{Abstract}
    We address emotional capability loss when an expressive text-to-speech
    (TTS) model trained on a high-resource language is adapted to a
    low-resource language. Using IndexTTS2 as the expressive base model, we
    propose an RL-enhanced adaptation framework for preserving emotional
    controllability, prosodic richness, and speech naturalness during transfer
    from Chinese to Taiwanese. The goal is to retain the pretrained emotional
    prior and realign it to limited target-language data.
  \end{block}

% ----------------------------------
% Motivation and objectives
% ----------------------------------
  \begin{alertblock}{Motivation and Objectives}
    Expressive low-resource TTS is limited by a tension between pronunciation
    adaptation and emotion preservation. Naive fine-tuning can improve
    intelligibility while flattening rhythm, pauses, emphasis, and intensity.

    \begin{itemize}
      \item \textbf{Emotional forgetting:} target-language fine-tuning can erase
      expressive behavior learned from high-resource speech.
      \item \textbf{Prosodic control loss:} adapted speech may stay intelligible
      but lose natural emphasis and intensity control.
      \item \textbf{Cross-lingual alignment:} emotional cues do not transfer
      perfectly across languages with different phonology and rhythm.
      \item \textbf{Objective:} recover controllable emotional expression in
      Taiwanese TTS through RL-based realignment.
    \end{itemize}
  \end{alertblock}

% ----------------------------------
% Method overview
% ----------------------------------
  \begin{block}{Proposed Framework}
    The method first adapts the expressive base model to Taiwanese
    pronunciation, then applies RL to recover controllable emotion and prosody.

    \begin{figure}
    \centering
    \begin{adjustbox}{width=0.94\linewidth,max height=0.26\textheight, center}
        \begin{tikzpicture}[
            node distance=1.25cm,
            posterstep/.style={rectangle, rounded corners=8pt, draw=black, very thick, fill=white, align=center, text width=12.2cm, minimum height=2.0cm},
            reward/.style={rectangle, draw=black, very thick, fill=white, align=center, text width=8.8cm, minimum height=1.7cm},
            connector/.style={ultra thick,->,>=stealth}
        ]
            \node[posterstep] (base) {\textbf{High-resource expressive IndexTTS2}\\Chinese expressive prior, speaker timbre, controllable emotion};
            \node[posterstep, below=of base] (adapt) {\textbf{Parameter-efficient Taiwanese adaptation}\\LoRA/adapters learn target-language pronunciation while limiting parameter drift};
            \node[posterstep, below=of adapt] (policy) {\textbf{GRPO emotional realignment}\\Text-to-semantic policy samples candidate semantic token sequences};
            \node[posterstep, below=of policy] (tts) {\textbf{Emotion-preserving Taiwanese TTS}\\Decoded waveform keeps content accuracy, timbre, intensity, and emphasis};

            \node[reward, right=2.4cm of policy, yshift=1.1cm] (emo) {Emotion category\\and intensity rewards};
            \node[reward, right=2.4cm of policy, yshift=-1.1cm] (prosody) {Local emphasis\\and ASR rewards};

            \draw[connector] (base) -- (adapt);
            \draw[connector] (adapt) -- (policy);
            \draw[connector] (policy) -- (tts);
            \draw[connector] (emo.west) -- (policy.east);
            \draw[connector] (prosody.west) -- (policy.east);
        \end{tikzpicture}
    \end{adjustbox}
    \caption{Two-stage adaptation and RL realignment pipeline.}
    \label{fig:method-flow}
\end{figure}


    \begin{itemize}
      \item \textbf{Stage 1:} LoRA or adapter-based adaptation limits parameter
      drift and reduces catastrophic forgetting.
      \item \textbf{Stage 2:} GRPO optimizes the autoregressive
      text-to-semantic policy with waveform-level rewards.
      \item \textbf{Anchor:} a KL penalty preserves pronunciation, timbre, and
      the adapted reference behavior.
    \end{itemize}
  \end{block}

% ----------------------------------
% Formulation
% ----------------------------------
  \begin{block}{RL Formulation}
    Semantic speech generation is formulated as an MDP over the IndexTTS2
    text-to-semantic module.

    \begin{itemize}
      \item \textbf{State:} input text, speaker timbre, duration constraint,
      target emotion, and prior semantic tokens.
      \item \textbf{Action:} the next discrete semantic speech token from
      policy $\pi_\theta$.
      \item \textbf{Reward:} a sequence-level score computed from the decoded
      waveform.
    \end{itemize}

    \[
      R = \lambda_1 R_{\mathrm{emo}}
        + \lambda_2 R_{\mathrm{intensity}}
        + \lambda_3 R_{\mathrm{emphasis}}
        + \lambda_4 R_{\mathrm{ASR}}
    \]

    \[
      \max_\theta\; \mathbb{E}[R(s,a)]
      - \beta\,\operatorname{KL}\!\left(\pi_\theta(\cdot|s)\,\|\,p_{\mathrm{adapt}}(\cdot|s)\right)
    \]
  \end{block}

% ----------------------------------
% Datasets and baselines
% ----------------------------------
  \begin{block}{Datasets and Baselines}
    The experiment is designed to separate basic language adaptation from
    emotional recovery, so each model is evaluated before adaptation, after
    low-resource adaptation, and after RL realignment.

    \begin{itemize}
      \item \textbf{Emotion supervision:} ESD and Expresso provide emotion
      labels, intensity cues, and emphasis annotations.
      \item \textbf{Taiwanese adaptation:} TAT\_MOE supports target-language
      adaptation and evaluation with low-resource subsets of roughly 1--5 hours.
      \item \textbf{Baselines:} high-resource IndexTTS2, naive fine-tuning,
      parameter-efficient adaptation, our RL-enhanced model, and CosyVoice3.
      \item \textbf{Core comparison:} Emotion Retention Gap and recovery ratio
      measure how much expressiveness is lost and then restored.
    \end{itemize}
  \end{block}

\end{column}

\separatorcolumn

\begin{column}{\colwidth}

% -------------------------------
% Architecture and reward design
% -------------------------------
  \begin{block}{Architecture and Reward Design}
    \begin{figure}
    \centering
    \includegraphics[width=0.95\linewidth,height=0.30\textheight,keepaspectratio]{images/architecture.pdf}
    \caption{RL-enhanced IndexTTS2 architecture with waveform-level reward models.}
    \label{fig:architecture}
\end{figure}


    Reward models evaluate decoded speech rather than only token likelihood.

    \begin{itemize}
      \item \textbf{Emotion category:} emotion2vec checks target emotion match.
      \item \textbf{Global intensity:} VAD distance from a neutral centroid
      encourages stronger affect.
      \item \textbf{Local emphasis:} pitch and energy patterns measure stressed
      words and expressive prosody.
      \item \textbf{Content accuracy:} ASR transcription compares generated
      speech with the input text.
    \end{itemize}
  \end{block}

% -------------------------------
% Evaluation
% -------------------------------
  \begin{block}{Evaluation Protocol}
    We evaluate intelligibility, speaker fidelity, emotional expressiveness,
    and prosodic control across three stages: source model, adapted model, and
    RL-optimized model.

    \begin{itemize}
      \item \textbf{Objective metrics:} WER from ASR, speaker similarity,
      emotion similarity or accuracy, Emotion Retention Gap, and recovery ratio.
      \item \textbf{Subjective MOS:} QMOS, SMOS, PMOS, and EMOS on a 1--5 scale
      for quality, speaker identity, prosody, and emotion match.
      \item \textbf{Fine-grained tests:} Emotion Intensity Test compares weak,
      medium, and strong generations; Emphasis Accuracy Test checks perceived
      stress locations.
      \item \textbf{Trade-off check:} gains in emotion should not degrade WER,
      speaker similarity, or naturalness.
    \end{itemize}
  \end{block}

% -------------------------------
% Preliminary results
% -------------------------------
  \begin{block}{Preliminary Taiwanese TTS Baseline}
    A pilot MOS evaluation with 12 SARC laboratory members used 20 sentences for
    intelligibility and naturalness plus 10 sentences for voice similarity,
    yielding 600 subjective rating samples.

    \begin{table}
    \centering
    \caption{Corpus-wide final report metrics. SER and SS are higher-is-better; WER is lower-is-better.}
    \renewcommand{\arraystretch}{1.12}
    \begin{tabular*}{\linewidth}{@{\extracolsep{\fill}}lccc}
      \toprule
      \textbf{Model} & \textbf{SER $\uparrow$} & \textbf{WER $\downarrow$} & \textbf{SS $\uparrow$} \\
      \midrule
      CosyVoice2 & 0.14 & 1.666$^{\dagger}$ & \textbf{0.581} \\
      CosyVoice3 & 0.13 & 2.110$^{\dagger}$ & 0.573 \\
      \midrule
      IndexTTS2-FT & 0.27 & \textbf{0.705} & 0.456 \\
      IndexTTS2-FT+VA/stress & 0.27 & 0.741 & 0.463 \\
      IndexTTS2-FT+RL & 0.28 & 0.816 & 0.433 \\
      \textbf{IndexTTS2-proposed} & \textbf{0.30} & 0.840 & 0.416 \\
      \bottomrule
    \end{tabular*}
    \vspace{0.25em}
    {\small $^{\dagger}$CosyVoice WER is a tone-sensitive diagnostic score and is not directly comparable to the IndexTTS2 rows.}
    \label{tab:final-report-metrics}
\end{table}


    CosyVoice2 reaches a mean MOS of 4.31, motivating a shift from basic quality
    toward emotional preservation and recovery.
  \end{block}

% -------------------------------
% Contributions
% -------------------------------
  \begin{exampleblock}{Expected Contributions}
    \begin{itemize}
      \item A low-resource emotional TTS adaptation framework separating
      pronunciation transfer from expressive realignment.
      \item A multi-objective reward design for emotion, intensity, emphasis,
      and semantic intelligibility.
      \item An evaluation protocol for emotional forgetting, recovery ratio,
      speaker similarity, and WER trade-offs.
    \end{itemize}
  \end{exampleblock}

% -------------------------------
% Team and references
% -------------------------------
  \begin{block}{Team and References}
    \scriptsize
    \textbf{Team:}
    Hao-Chun Hsieh -- formulation and methodology;
    Pao-Hua Hsu -- documentation, visualization, poster, and presentation;
    Nien-Tse Wu -- data, evaluation, implementation, and tuning.

    \vspace{0.4em}
    \textbf{References:}
    Zhou et al. IndexTTS2, 2025.\quad
    Li et al. EMORL-TTS, 2025.\quad
    Ma et al. emotion2vec, 2023.\quad
    Liao. TAT\_MOE Corpus, 2022.\quad
    Du et al. CosyVoice2, 2024.
  \end{block}

\end{column}
\separatorcolumn

\end{columns}
\end{frame}

\end{document}
